\begin{document}
\maketitle

\begin{abstract}
Let $\A\subset\B$ be an irreducible finite-index inclusion of conformal
nets on $S^1$, let $\A=\B^{\Z_4}$ for a proper gauge action commuting
with rotations, and suppose that $\B$ is holomorphic, with
$\mu_{\B}=1$ and central charge $8$.  We prove the strongest correct
form of the proposed KMS extension statement.  At angular inverse
temperature $\beta>0$, the vacuum-sector Gibbs state of $\B$ restricts
to a rotational KMS state $\omega_\beta$ of $\A$.  The group average
\[
 E=\frac14\sum_{r=0}^{3}\alpha_r
\]
is the unique normal faithful conditional expectation
$\B(I)\to\A(I)$ on every interval, and the extension is
$\widetilde\omega_\beta=\omega_\beta\circ E$.  It is the unique
$\Z_4$-invariant state extension, even before the KMS condition is
imposed.  Finally, the Kawahigashi--Longo--M\"uger sector theorem
together with the rotational KMS classification of
Longo--Tanimoto implies that a holomorphic conformal net has only one
rotational $\beta$-KMS state.  Hence every rotational KMS extension is
automatically $\Z_4$-invariant and equals
$\omega_\beta\circ E$.

Two qualifications are essential.  The vacuum \emph{vector} state is
a rotation ground state, not a finite-temperature rotational KMS
state; ``vacuum KMS'' below means the vacuum-\emph{sector Gibbs
state}.  Also, $\mu=1$ eliminates DHR sectors but does not by itself
classify thermal states: the Longo--Tanimoto theorem is the necessary
analytic bridge.  We quote in full the precise forms of Takesaki's
expectation theorem, the finite-index/Q-system theorem, the KLM
$\mu$-index theorem, Araki's CAR KMS theorem, and the rotational KMS
classification, and verify their hypotheses against the frozen TFPT
finite-collar data.  The final section isolates the remaining
continuum identification in one box.
\end{abstract}

\tableofcontents

\section{Setting and correction of terminology}

\subsection{Conformal-net hypotheses}

Let $I\mapsto\B(I)$ be an irreducible local conformal net on $S^1$ in
its vacuum representation $(\Hh_{\B},U_{\B},\Omega_{\B})$.  Let
$G=\Z_4$ act properly by vacuum-preserving internal symmetries
$\alpha_r$, meaning
\[
 \alpha_r(\B(I))=\B(I),\qquad
 \alpha_r\Ad U_{\B}(g)=\Ad U_{\B}(g)\alpha_r
\]
for every interval $I$, every conformal transformation $g$, and
$r\in\Z_4$.  Put
\[
 \A(I)=\B(I)^G.
\]
We assume that the local inclusion is irreducible and has Jones index
four:
\[
 [\B(I):\A(I)]=4,\qquad
 \A(I)'\cap\B(I)=\C\id.
\tag{1.1}\label{eq:local-inclusion}
\]
For a proper outer action of a finite group on a factor, the first
identity follows from the fixed-point theorem; it is also included
explicitly among our hypotheses so no outerness convention is hidden.

The net $\B$ is assumed holomorphic in the operator-algebraic sense:
it is completely rational and
\[
 \mu_{\B}=1.
\tag{1.2}\label{eq:holomorphic}
\]
The central charge $c=8$ identifies the intended endpoint as the
$(E_8)_1$ net, but is not needed for KMS uniqueness once
\eqref{eq:holomorphic} is known.

\subsection{The rotational KMS state}

Let $L_0^{\B}\geq0$ be the conformal Hamiltonian and let
$\tau_t=\Ad e^{itL_0^{\B}}$ be the rotation dynamics on the locally
normal universal algebra $C^*_{\mathrm{ln}}(\B)$.  For $\beta>0$ put
\[
 \Phi_\beta^{\B}(x)
 =\frac{\Tr_{\Hh_{\B}}(e^{-\beta L_0^{\B}}x)}
 {\Tr_{\Hh_{\B}}(e^{-\beta L_0^{\B}})}.
\tag{1.3}\label{eq:gibbs}
\]
For a completely rational conformal net the character is finite, and
the rotational KMS classification quoted below in fact recovers the
trace-class property from any extremal KMS state.  In the concrete
$(E_8)_1$ case it is also immediate from the character
$E_4/\eta^8$.

\begin{definition}[Geometric rotational state used here]
\label{def:geometric}
The state on the subnet is
\[
 \omega_\beta:=\Phi_\beta^{\B}|_{C^*_{\mathrm{ln}}(\A)}.
\tag{1.4}\label{eq:omega-def}
\]
At the canonical angular normalization we set
$\beta=\beta_{\mathrm{angle}}=2\pi$.
\end{definition}

\begin{remark}[Ground state versus rotational KMS state]
The vector functional
$\langle\Omega_{\B},\,\cdot\,\Omega_{\B}\rangle$ is a ground state for
rotations.  It is not the finite-temperature state
\eqref{eq:gibbs}.  The geometric KMS state on the real line for
translations/dilations and the rotational Gibbs state on the circle
are likewise different constructions.  The theorem below concerns
the latter, because its unconditional classification by
superselection sectors is available in exactly this setting.
\end{remark}

\section{Classical theorems used, stated in full}

\subsection{Araki's CAR KMS theorem}

\begin{theorem}[Araki: unique quasi-free CAR equilibrium state]
\label{thm:araki}
Let $\mathfrak h$ be a complex Hilbert space, let $h=h^*$ be a
possibly unbounded one-particle Hamiltonian, and let
$\tau_t^h$ be the strongly continuous Bogoliubov dynamics on
$\CAR(\mathfrak h)$ determined on generators by
\[
 \tau_t^h(a(f))=a(e^{ith}f).
\]
For every $\beta>0$ there is a gauge-invariant quasi-free state
$\varphi_{\beta,h}$ whose two-point function is
\[
 \varphi_{\beta,h}(a^*(f)a(g))
 =\langle g,C_{\beta,h}f\rangle,\qquad
 C_{\beta,h}=(\id+e^{\beta h})^{-1},
\tag{2.1}\label{eq:araki-cov}
\]
and whose higher even correlation functions are the determinants
prescribed by the CAR Wick rule, while odd correlation functions
vanish.  This state is a $\beta$-KMS state for $\tau^h$.  In the
factorial self-dual CAR representation (in particular when the
covariance has no $0$ or $1$ eigenvalue in the finite-dimensional
case), it is the unique $\beta$-KMS state for that Bogoliubov
dynamics.  Thus the KMS state is completely determined by its
two-point covariance.
\end{theorem}

This is the KMS uniqueness theorem in
H.~Araki, \emph{On quasifree states of CAR and Bogoliubov
automorphisms}, Publ.\ Res.\ Inst.\ Math.\ Sci.\ \textbf{6}
(1970/71), 385--442, Theorem~3 and Eq.~(6.9), translated from the
self-dual notation to the ordinary gauge-invariant CAR notation.
Araki's factor proviso is recorded rather than suppressed.

\subsection{Takesaki's conditional-expectation theorem}

\begin{theorem}[Takesaki: modular invariance criterion]
\label{thm:takesaki}
Let $M$ be a von Neumann algebra, let $\phi$ be a faithful normal
semifinite weight on $M$, and let $N\subset M$ be a von Neumann
subalgebra such that $\phi|_N$ is semifinite.  Write
$\sigma^\phi=(\sigma_t^\phi)_{t\in\R}$ for the modular automorphism
group of $\phi$.  The following are equivalent:
\begin{enumerate}[label=\textup{(\roman*)}]
\item $N$ is globally invariant under the modular group,
  $\sigma_t^\phi(N)=N$ for every $t\in\R$;
\item there exists a normal faithful conditional expectation
  $E_\phi:M\to N$ such that
  $\phi\circ E_\phi=\phi$ on $M_+$.
\end{enumerate}
When these conditions hold, $E_\phi$ is unique among the normal
conditional expectations preserving $\phi$, it commutes with the
modular group,
\[
 E_\phi\sigma_t^\phi=\sigma_t^{\phi|_N}E_\phi,
\tag{2.2}\label{eq:takesaki-intertwine}
\]
and the modular group of the restricted weight is the restriction:
\[
 \sigma_t^{\phi|_N}=\sigma_t^\phi|_N.
\tag{2.3}\label{eq:mod-restrict}
\]
\end{theorem}

This is M.~Takesaki, \emph{Conditional expectations in von Neumann
algebras}, J.\ Funct.\ Anal.\ \textbf{9} (1972), 306--321, and
\emph{Theory of Operator Algebras II}, Springer, 2003,
Theorem~IX.4.2.

\subsection{Finite index, Radon--Nikodym uniqueness, and Q-systems}

\begin{theorem}[Finite-index expectation and Q-system package]
\label{thm:qsystem}
Let $N\subset M$ be an inclusion of infinite factors with finite
index and separable predual.
\begin{enumerate}[label=\textup{(\alph*)}]
\item There is a normal faithful conditional expectation
  $E_0:M\to N$ of finite index.  Fix one such expectation.  Every
  other normal faithful conditional expectation $F:M\to N$ is of
  the form
  \[
    F(x)=E_0(h^{1/2}xh^{1/2}),\qquad x\in M,
  \tag{2.4}\label{eq:rn-expectation}
  \]
  for a positive invertible $h\in N'\cap M$ satisfying
  $E_0(h)=\id$.  Conversely every such $h$ defines an expectation.
  In particular, if $N'\cap M=\C\id$, then the normal faithful
  conditional expectation is unique.
\item Let $\iota:N\hookrightarrow M$ be the inclusion morphism and
  $\bar\iota:M\to N$ a conjugate.  The dual canonical endomorphism
  $\theta=\bar\iota\iota\in\operatorname{End}(N)$, together with
  standard conjugacy intertwiners
  \[
    w\in\Hom(\id_N,\theta),\qquad
    x\in\Hom(\theta,\theta^2),
  \]
  is a standard Q-system: $w$ is a unit, $x$ is associative and
  unital with $w$, and $x$ obeys the Frobenius and standardness
  relations.  Conversely, every standard Q-system in
  $\operatorname{End}(N)$ reconstructs, up to isomorphism fixing
  $N$, a finite-index extension $N\subset M$.  Its dimension is
  \[
    d(\theta)=[M:N].
  \tag{2.5}\label{eq:q-index}
  \]
  If the endomorphisms lie in the braided DHR category of a
  conformal net, the reconstructed extension is local exactly when
  the Q-system is commutative with respect to the braiding.
\item If $M$ carries an outer action of a finite group $G$ and
  $N=M^G$, then
  \[
    E_0(x)=\frac1{|G|}\sum_{g\in G}\alpha_g(x),\qquad
    [M:N]=|G|.
  \tag{2.6}\label{eq:group-average}
  \]
  The dual canonical endomorphism is the regular Q-system.  For
  $G=\Z_4$ it is the sum of the four one-dimensional character
  sectors and has dimension four.
\end{enumerate}
\end{theorem}

Part (a) is the Radon--Nikodym description of expectations in the
finite-index theory of H.~Kosaki and R.~Longo; see R.~Longo,
\emph{Index of subfactors and statistics of quantum fields I},
Commun.\ Math.\ Phys.\ \textbf{126} (1989), 217--247, especially
the uniqueness/minimal-index discussion, and
F.~Hiai, \emph{Minimum index for subfactors and entropy II},
J.\ Math.\ Soc.\ Japan \textbf{43} (1991), 347--379.  Parts (b) and
(c) are the subfactor/Q-system and fixed-point forms used in
R.~Longo and K.-H.~Rehren, \emph{Nets of subfactors},
Rev.\ Math.\ Phys.\ \textbf{7} (1995), 567--597, together with the
local-extension formulation of J.~B\"ockenhauer and D.~E.~Evans,
\emph{Modular invariants, graphs and $\alpha$-induction for nets of
subfactors I}, Commun.\ Math.\ Phys.\ \textbf{197} (1998), 361--386.

\subsection{\texorpdfstring{The KLM $\mu$-index theorem}{The KLM mu-index theorem}}

\begin{theorem}[Kawahigashi--Longo--M\"uger]
\label{thm:klm}
Let $\mathcal N$ be a local irreducible conformal net satisfying the
split property and strong additivity, and suppose that its
two-interval inclusion
\[
 \mathcal N(E)\subset\widehat{\mathcal N}(E)
 :=\mathcal N(E')'
\]
has finite index.  Then $\mathcal N$ is completely rational; it has
finitely many inequivalent irreducible DHR sectors
$\rho_0=\id,\rho_1,\ldots,\rho_n$, all of finite statistical
dimension, every locally normal representation is a direct sum of
these sectors, the braiding is non-degenerate, and
\[
 \mu_{\mathcal N}
 =[\widehat{\mathcal N}(E):\mathcal N(E)]
 =\sum_{j=0}^{n}d(\rho_j)^2.
\tag{2.7}\label{eq:mu-sum}
\]
The two-interval inclusion is the Longo--Rehren inclusion associated
with the full sector category.  If
$\mathcal N\subset\mathcal M$ is a finite-index irreducible local
extension of completely rational nets, then
\[
 \mu_{\mathcal M}
 =\frac{\mu_{\mathcal N}}{[\mathcal M:\mathcal N]^2}.
\tag{2.8}\label{eq:mu-extension}
\]
Consequently, $\mu_{\mathcal M}=1$ if and only if the vacuum is the
only irreducible DHR sector of $\mathcal M$.
\end{theorem}

This is Y.~Kawahigashi, R.~Longo and M.~M\"uger,
\emph{Multi-interval subfactors and modularity of representations in
conformal field theory}, Commun.\ Math.\ Phys.\ \textbf{219} (2001),
631--669: the two-interval/Longo--Rehren identification, the equality
of $\mu$ with the global sector dimension, and the extension-index
formula are the results used here.  The final equivalence follows
from \eqref{eq:mu-sum}, since $d(\rho_0)=1$ and every nonzero
statistical dimension is at least one.

\subsection{Rotational KMS classification}

\begin{theorem}[Longo--Tanimoto]
\label{thm:lt}
Let $\mathcal N$ be a completely rational conformal net and let
$C^*_{\mathrm{ln}}(\mathcal N)$ be its locally normal universal
$C^*$-algebra.  Then
\[
 C^*_{\mathrm{ln}}(\mathcal N)
 \cong\bigoplus_{j=0}^{n}B(\Hh_j),
\tag{2.9}\label{eq:universal-sum}
\]
where the summands correspond bijectively to the inequivalent
irreducible DHR representations $\rho_j$.  Every extremal rotational
$\beta$-KMS state is the Gibbs state in one irreducible
representation:
\[
 \Phi_{j,\beta}(x)
 =\frac{\Tr_{\Hh_j}(e^{-\beta L_0^{(j)}}\rho_j(x))}
 {\Tr_{\Hh_j}(e^{-\beta L_0^{(j)}})}.
\tag{2.10}\label{eq:sector-gibbs}
\]
In particular $e^{-\beta L_0^{(j)}}$ is trace class for every sector
that supports such a state.  Every rotational $\beta$-KMS state is a
convex combination of the extremal sector Gibbs states.  Hence, if
the vacuum is the only irreducible sector, the vacuum-sector Gibbs
state is the unique rotational $\beta$-KMS state.
\end{theorem}

Equation \eqref{eq:universal-sum} is the completely rational universal
algebra theorem of Carpi--Conti--Hillier--Weiner used in
R.~Longo and Y.~Tanimoto, \emph{Rotational KMS states and type I
conformal nets}, Commun.\ Math.\ Phys.\ \textbf{357} (2018),
249--266, Theorems~3.1, 3.2 and 3.7.  The last sentence is an
immediate specialization to a one-summand algebra.

\section{Verification of the hypotheses}

\subsection{The index-four fixed-point Q-system}

For each interval set $N=\A(I)$ and $M=\B(I)$.  Conformal-net local
algebras are factors; \eqref{eq:local-inclusion} gives finite index
and trivial relative commutant.  Theorem~\ref{thm:qsystem} therefore
applies.  The group average is
\[
 E_I(x)=\frac14\bigl(x+\alpha_1(x)+\alpha_2(x)+\alpha_3(x)\bigr).
\tag{3.1}\label{eq:EI}
\]
It is normal, unital, positive, faithful, idempotent, and
$\A(I)$-bimodular.  For example, if $a,b\in\A(I)$, then
\[
 E_I(axb)=\frac14\sum_{r=0}^{3}a\alpha_r(x)b
 =aE_I(x)b.
\tag{3.2}\label{eq:bimodule}
\]
Faithfulness follows because $x\geq0$ and $E_I(x)=0$ imply that the
sum of four positive operators is zero, hence $x=0$.  Compatibility
under interval inclusion is immediate from the same formula, so the
$E_I$ define a net expectation and a global group average
\[
 E:C^*_{\mathrm{ln}}(\B)\longrightarrow
 C^*_{\mathrm{ln}}(\A).
\tag{3.3}\label{eq:global-E}
\]

The regular $\Z_4$ Q-system has
\[
 \theta\simeq\rho_0\oplus\rho_1\oplus\rho_2\oplus\rho_3,
\qquad d(\rho_r)=1,\qquad d(\theta)=4.
\tag{3.4}\label{eq:regular-q}
\]
The frozen TFPT simple-current arithmetic
$h(J^r)=1$ for $r=1,2,3$ verifies the integer-spin/locality condition
for the intended index-four extension.  The arithmetic cited in
\texttt{v973} is
\[
 \mu_{\A}=16,\qquad [\B:\A]=4,\qquad
 \mu_{\B}=\frac{16}{4^2}=1.
\tag{3.5}\label{eq:tfpt-mu}
\]
This is precisely the KLM extension formula
\eqref{eq:mu-extension}; the Python computation is a finite exact
check of its numerical inputs, not a replacement for the theorem.

